\documentclass[letterpaper,10pt]{texMemo}

\usepackage{graphicx}

\memoto{Dr. Mark Yoder}
\memofrom{Michael McDonald}
\memosubject{Week 6 Memo}
\memodate{\today}
\logo{\includegraphics[width=0.5\textwidth]{RoseLogo.png}}

\begin{document}
\maketitle

\section{Homework 6 Summary}
\noindent Definitions for CONFIG\_CPU\_V6, CONFIG\_CPU\_V6K, and CONFIG\_CPU\_V7
can be found in \\
\verb .../arch/arm/mm/Kconfig  and look something like this:
\begin{verbatim}
 # ARMv7
config CPU_V7
	bool "Support ARM V7 processor" if ARCH_INTEGRATOR || MACH_REALVIEW_EB
|| MACH_REALVIEW_PBX
	select CPU_32v6K
	select CPU_32v7
	select CPU_ABRT_EV7
	select CPU_CACHE_V7
	select CPU_CACHE_VIPT
	select CPU_COPY_V6 if MMU
	select CPU_CP15_MMU if MMU
	select CPU_CP15_MPU if !MMU
	select CPU_HAS_ASID if MMU
	select CPU_PABRT_V7
	select CPU_TLB_V7 if MMU
\end{verbatim}
\noindent This basically creates a dependency tree of necessary components that
must be initialized for this processor. The actual selection of which
definition we use occurs in \verb KERNEL/.config , which looks like this:

\begin{verbatim}
#
# Processor Type
#
CONFIG_CPU_V7=y
CONFIG_CPU_32v6K=y
CONFIG_CPU_32v7=y
CONFIG_CPU_ABRT_EV7=y
CONFIG_CPU_PABRT_V7=y
CONFIG_CPU_CACHE_V7=y
CONFIG_CPU_CACHE_VIPT=y
CONFIG_CPU_COPY_V6=y
CONFIG_CPU_TLB_V7=y
CONFIG_CPU_HAS_ASID=y
CONFIG_CPU_CP15=y
CONFIG_CPU_CP15_MMU=y
\end{verbatim}

\newpage
\noindent The command \verb b  \verb start_kernel  can be found at line 104 in
the file \verb KERNEL/arch/arm/kernel/head-common.S ,\\ which is included at the
bottom of \verb KERNEL/arch/arm/kernel/head.S . This calls the
\verb start_kernel()  function in main.c, which is located in
\verb KERNEL/init/main.c , located at line 473.


\section{Exercise Summary}

\begin{center}
\begin{tabular}{|p{4cm}|l|l|p{8cm}|}
\hline
\textbf{Objective} & \textbf{Exercise} &\textbf{Status} & \textbf{Notes}\\
\hline
 Boot sequence & Exercise 21a & Completed & The FAT partition seemed to
already be loaded in \verb /media/BEAGLEBONE \\
\hline
 Systemd & Exercise 21b & Completed & The exercise reminded me that I need to
learn JavaScript.\\
\hline
\end{tabular}
\end{center}


\noindent The main points I took away from these exercises were the following:
\begin{itemize}
 \item The kernel is incredibly complex, and contains many files that have the
same name, but are located in different paths, which leads to some serious
confusion. 
\end{itemize}

\end{document}
